\section{Phase markers}
\label{sec:markers}

\subsection{Structure factor and the AFM label}
\label{sec:sq}

To identify the antiferromagnetic phase we use the structure factor
\begin{equation}
S(q)=\frac{1}{L}\sum_{i,j}e^{iq(i-j)}\langle Z_iZ_j\rangle,
\end{equation}
where $\langle Z_iZ_j\rangle$ is the ground-state correlator and $L=8$
throughout this section.
That is, $S(q)$ is the Fourier weight of N\'eel-type correlations, so
long-range antiferromagnetic order appears as a Bragg-like peak at
$q=\pi$. The $S(q)$ curves at one representative point in each
of the three phases (labels trivial / topo / AFM, coordinates in the
dataset) on the shared $q$ grid ($99$ points including $q=\pi$) are shown in
Fig.~\ref{fig:D03}.
The $q=\pi$ component is used as the AFM diagnostic throughout the
hardware-facing protocol. Its empirical phase discrimination is reported
with the results rather than assumed as part of the measurement method.

\begin{figure}[htbp]
\centering
\includegraphics[width=\columnwidth]{exp02_D03.pdf}
\caption{Structure factor $S(q)$ at three phase-representative points (one
per phase) on the shared $q$ grid; only the AFM point peaks at $q=\pi$.
\label{fig:D03}}
\end{figure}
